% Answers to Chapter 1, translated from sv/back/facit/01.tex.

\facitavsnitt{c1}{Chapter 1}{Number systems and basic arithmetic}
\begin{facit}

\svar{1.1} \sv{a} $13$ \sv{b} $25$ \sv{c} $5$ \sv{d} $15$

\svar{1.2} \sv{a} $\frac{5}{12}$ \sv{b} $\frac{7}{20}$ \sv{c} $\frac{3}{2} = 1.5$ \sv{d} $2$

\svar{1.3} \sv{a} $100$ \sv{b} $25\,\%$ \sv{c} Lowest $95\,\Omega$, highest $105\,\Omega$.

\svar{2.1} \sv{a} $\frac{1}{4}\,\Omega^{-1}$ \sv{b} $4\,\Omega$

\svar{2.2} \sv{a} $R_{\text{p}} = 6\,\text{k}\Omega$ with both formulas.
\sv{b} $10\,\text{k}\Omega$
\sv{c} The parallel connection works like a bracket: its value must be known before $R_1$ can be
added. Adding all three resistors gives the wrong answer, $36\,\text{k}\Omega$.

\svar{2.3} \sv{a} $2\,\text{mA}$ \sv{b} $8\,\text{V}$ \sv{c} $12\,\text{V}$
\sv{d} $I_2 = 0.5\,\text{mA}$, $I_3 = 1.5\,\text{mA}$; the sum is $2\,\text{mA} = I$.

\svar{2.4} \sv{a} $8\,\text{V}$ \sv{b} $5.3\,\text{V}$

\svar{2.5} \sv{a} $75\,\%$ \sv{b} $21.6\,\text{W}$

\svar{2.6} \sv{a} $\mathbb{N}$ \sv{b} $\mathbb{Z}$ \sv{c} $\mathbb{Q}$ \sv{d} Irrational
\sv{e} $\mathbb{Q}$ \sv{f} Irrational

\svar{3.1} \sv{a} $-13$ \sv{b} $0$ \sv{c} $-6$ \sv{d} $16$ \sv{e} $1$

\svar{3.2} \sv{a} $-\abs{4} < -2.5 < 0 < \abs{-3} < \frac{7}{2}$ \sv{b} $4$
\sv{c} Both are $5$: the absolute value of a difference is the distance between the numbers on the
number line.
\sv{d} $\abs{U_1 - U_2} = 0.3\,\text{V}$

\svar{3.3} $\frac{1}{4} = 0.25 = 25\,\%$;\enspace $\frac{3}{5} = 0.6 = 60\,\%$;\enspace
$\frac{1}{8} = 0.125 = 12.5\,\%$;\enspace $\frac{3}{8} = 0.375 = 37.5\,\%$

\svar{3.4} \sv{a} $\frac{1}{3}$ \sv{b} $1$ \sv{c} $1$ \sv{d} $2$
\sv{e} The parallel resistance of two $4\,\Omega$ resistors, which is $2\,\Omega$.

\svar{3.5} \sv{a} $25\,\%$ \sv{b} $15\,\%$
\sv{c} No. $1.10 \times 0.90 = 0.99$, so the value ends up $1\,\%$ lower than at the start.
\sv{d} $75\,\text{W}$

\svar{3.6} \sv{a} $4.23\,\text{k}\Omega$ to $5.17\,\text{k}\Omega$ \sv{b} Yes.
\sv{c} About $4.3\,\%$
\sv{d} No. The permitted range is $209$--$231\,\Omega$, and $208\,\Omega$ deviates by about
$5.5\,\%$.

\svar{3.7} \sv{a} True: $n = \frac{n}{1}$.
\sv{b} False, for example $\frac{1}{2}$.
\sv{c} False: $\sqrt{9} = 3$.
\sv{d} False, for example $\sqrt{2} + (-\sqrt{2}) = 0$.
\sv{e} True: $0.333\ldots = \frac{1}{3}$.
\sv{f} False, for example $\pi$ and $\sqrt{2}$.

\svar{3.8} \sv{a} $50\,\Omega$ \sv{b} $25\,\Omega$
\sv{c} $n$ equal terms give $\frac{1}{R_{\text{p}}} = \frac{n}{R}$.
\sv{d} $5$ resistors

\svar{3.9} \sv{a} $20\,\Omega$ \sv{b} $280\,\Omega$ \sv{c} Yes, $12\,\Omega < 20\,\Omega$.
\sv{d} No. A parallel connection is always smaller than the smallest of the resistors in it, here
$30\,\Omega$.

\svar{3.10} From the top: $I = 3\,\text{mA}$;\enspace $U = 11\,\text{V}$;\enspace
$R = 3\,\text{k}\Omega$;\enspace $I = 5\,\text{mA}$

\svar{3.11} \sv{a} $190\,\text{W}$ \sv{b} $152\,\text{W}$ \sv{c} $76\,\%$ \sv{d} $48\,\text{W}$

\svar[Test yourself]{T} \sv{1} The irrational numbers, and so $\mathbb{R}$: $5$ is not the square
of an integer, so $\sqrt{5}$ cannot be written as a fraction.
\sv{2} $3\,\Omega$

\end{facit}
